\chapter*{Conclusion Générale}
\addcontentsline{toc}{chapter}{Conclusion Générale}

Ce projet de fin d'études m'a permis de mettre en pratique les connaissances acquises au cours de ma formation académique et de développer des compétences professionnelles essentielles dans le domaine du génie informatique.

\section*{Réalisations accomplies}

Au cours de ce stage, j'ai pu :

\begin{itemize}
    \item Concevoir et implémenter un système de synchronisation incrémentale (Delta Sync) robuste et performant
    \item Maîtriser l'architecture hexagonale et le pattern CQRS
    \item Développer en utilisant les technologies modernes (Symfony, ReactJS, MySQL)
    \item Collaborer efficacement dans une équipe agile utilisant SCRUM
    \item Mettre en place des tests et du contrôle de qualité
    \item Assurer la documentation et la maintenabilité du code
end{itemize}

\section*{Apprentissages et compétences acquises}

Ce stage m'a enrichi sur plusieurs plans :

\begin{itemize}
    \item \textbf{Techniques} : Maîtrise approfondie de Symfony, ReactJS, et des architectures modernes
    \item \textbf{Méthodologiques} : Application pratique de SCRUM et des bonnes pratiques DevOps
    \item \textbf{Professionnels} : Communication, travail en équipe, gestion de projet
    \item \textbf{Personnels} : Autonomie, prise de responsabilité, persévérance
end{itemize}

\section*{Perspectives futures}

Les travaux futurs pourraient inclure :

\begin{itemize}
    \item Optimisation supplémentaire des performances du système de synchronisation
    \item Implémentation de mécanismes de cache distribué
    \item Extension du système à d'autres modules de PowerTeam
    \item Amélioration de l'interface utilisateur et de l'expérience utilisateur
end{itemize}
